\documentclass[12pt]{article}
\usepackage[utf8]{inputenc}
\usepackage[T1]{fontenc}
\usepackage{hyperref}
\title{Tay Boi Language Resources for K-12 Curriculum Development}
\author{Compiled by Claude Sonnet 5 (claude-sonnet-5) for Andr\'as Kornai}
\date{}
\begin{document}
\maketitle

\section{Introduction}
Tay Boi (T\^ay B\h{\^o}i; ISO 639-3 \texttt{tas}, Glottocode tayb1240) is an extinct French-based pidgin of Vietnam. According to the search summary it arose in the 1860s around Saigon, was used mainly by Vietnamese people working in French households, and fell out of use after 1954, probably dying out by the 1980s. It appears never to have been learned as a first language. This report is based on a single web search and found only academic and reference sources.

\section{Primary Educational Resources}
No primers, textbooks, courses, or school programmes were found. There is no known speaker community or school context.

\section{Linguistic and Reference Materials}
\begin{itemize}
\item Susan Love (2000), University of Adelaide MA thesis on Tay Boi (full text PDF). \url{https://web4.library.adelaide.edu.au/dspace/bitstream/2440/101230/1/02whole.pdf}
\item Front matter of the same thesis. \url{https://web4.library.adelaide.edu.au/dspace/bitstream/2440/101230/2/01front.pdf}
\item Repository record for the thesis. \url{https://web4.library.adelaide.edu.au/dspace/handle/2440/101230?mode=full}
\item ``Emergent systematicity in T\^ay B\h{\^o}i (Vietnamese Pidgin French)'', 2020 conference poster at Cambridge Open Engage. \url{https://www.cambridge.org/engage/coe/article-details/5fa931acbf697b0012cceeed}
\item Wikipedia, ``T\^ay B\h{\^o}i Pidgin French'' (general overview, secondary source). \url{https://en.wikipedia.org/wiki/T%C3%A2y_B%E1%BB%93i_Pidgin_French}
\item Glottolog entry for Tay Boi, with bibliography including Reinecke (1971) and Love (2000). \url{https://glottolog.org/resource/languoid/id/tayb1240}
\end{itemize}

\section{Religious and Community Resources}
None found.

\section{Digital and Audio Resources}
No audio, video, or corpus resources were found.

\section{Curriculum Development Context}
There are no native speakers, so any use would be a heritage or history-of-language unit, for example within lessons on Vietnamese colonial history or on contact languages. The Love thesis is the likeliest source of attested word lists and sentences, but I did not inspect its contents and cannot say how much primary text it holds. A full revitalization curriculum is not realistic on the evidence found.

\end{document}